\documentclass[11pt]{article}
\usepackage{../../latex/stat4ml_style}

\title{Chapter 8: Probabilistic Sequence and Language Models \\ \large Part III --- Statistical Foundations of Deep Learning \& Reinforcement Learning}
\author{Omid Safarzadeh}
\date{}

\begin{document}
\maketitle

\begin{abstract}
\noindent
A language model is a probability distribution over sequences, and everything it
does follows from one identity: the chain rule of probability, which factorises
the joint probability of a sentence into a product of next-token conditionals.
Training such a model is maximum likelihood on that factorisation --- so its loss
is cross-entropy and its reported score, perplexity, is the exponentiated
cross-entropy of Chapter~5. This chapter builds the autoregressive model from the
chain rule, places $n$-gram and neural models on the same probabilistic footing,
and then reads two pieces of the modern stack through a statistical lens: word
embeddings, which turn out to factorise a pointwise-mutual-information matrix, and
attention, which is kernel-weighted averaging --- Nadaraya--Watson regression ---
in disguise.
\end{abstract}

\tableofcontents

% ======================================================================
\section{The chain rule and the autoregressive model}
% ======================================================================

\begin{theorem}[Chain rule of probability]
For any sequence $x_{1:n}=(x_1,\dots,x_n)$,
\[
  p(x_{1:n}) = \prod_{t=1}^n p(x_t\mid x_{1:t-1}).
\]
\end{theorem}
\begin{proof}
Repeated application of the definition of conditional probability gives
\[
  p(x_{1:n}) = p(x_{1:n-1})\,p(x_n\mid x_{1:n-1}),
\]
and unrolling the first factor the same way down to $p(x_1)$ gives the product.
\end{proof}

This exact, assumption-free factorisation is the \keyterm{autoregressive} model:
to model a sequence it suffices to model the next token given the past. An
\keyterm{$n$-gram} model makes the Markov approximation
\[
  p(x_t\mid x_{1:t-1}) \approx p(x_t\mid x_{t-n+1:t-1})
\]
--- condition on the last $n-1$ tokens only --- and estimates each conditional by
counting. A neural language model keeps the
exact conditioning on all of $x_{1:t-1}$ but parameterises $p(x_t\mid x_{1:t-1})$
with a network. Either way the object is the same conditional distribution.

Fitting it is maximum likelihood: maximise $\sum_t \log p(x_t\mid x_{1:t-1})$,
which is minimising cross-entropy against the data (Chapter~2). So the loss is the
cross-entropy of Chapter~5, and the model's \keyterm{perplexity} is
$e^{\text{cross-entropy}}$ --- the effective number of choices it is torn between
at each token. The left panel shows held-out perplexity dropping as the context
grows from unigram to trigram on a second-order source, then ticking up for the
$4$-gram as the counts grow sparse: more context helps only while the data can
still estimate it.

% ======================================================================
\section{Embeddings factorise a PMI matrix}
% ======================================================================

Before a token enters a network it is an \keyterm{embedding}, a dense vector, and
the surprising fact is what those vectors encode. \keyterm{word2vec}'s
skip-gram-with-negative-sampling objective, trained on which words co-occur, has
an exact interpretation.

\begin{theorem}[word2vec is implicit matrix factorisation]
Skip-gram with negative sampling (with $k$ negatives) is, at its optimum,
factorising the matrix of \keyterm{shifted pointwise mutual information}:
word and context vectors satisfy
\[
  v_w^\top c_c \;=\; \mathrm{PMI}(w,c) - \log k,
  \qquad \mathrm{PMI}(w,c)=\log\frac{p(w,c)}{p(w)\,p(c)} .
\]
\end{theorem}

The content is that semantics is geometry: PMI is large for words that co-occur
more than chance, so factorising it places words that share contexts near one
another. The right panel builds a PMI matrix from a toy two-topic corpus and
factorises it; the monarchy words and the animal words fall into two separated
clusters, with no labels ever supplied. This is the same low-rank,
SVD-of-an-association-matrix idea as PCA (Part~II) --- embeddings are PCA for
co-occurrence.

% ======================================================================
\section{Attention is kernel smoothing}
% ======================================================================

The transformer's one new operation is \keyterm{attention}: a query attends to a
set of key--value pairs and returns a weighted average of the values. Written out,
with query $q$, keys $k_i$, values $v_i$,
\[
  \operatorname{attn}(q) = \sum_i \alpha_i\, v_i,
  \qquad \alpha_i = \frac{\exp(q^\top k_i/\sqrt d)}{\sum_j \exp(q^\top k_j/\sqrt d)} .
\]
This is a statistical object a century old.

\begin{definition}[Nadaraya--Watson estimator]
The kernel-regression estimate of $\mathbb{E}[v\mid q]$ from pairs $(k_i,v_i)$ is
\[
  \hat m(q) = \frac{\sum_i K(q,k_i)\,v_i}{\sum_i K(q,k_i)},
\]
a weighted average of the values with weights given by a similarity kernel $K$.
\end{definition}

Attention \emph{is} the Nadaraya--Watson estimator with the kernel
$K(q,k)=\exp(q^\top k/\sqrt d)$: the softmax normalisation is exactly the
denominator $\sum_i K$, and the output is exactly the kernel-weighted mean of the
values. The middle panel makes this concrete --- attention over scattered
key--value points traces the same smooth regression curve a kernel smoother
would, with the softmax weights for a query peaking at the nearby keys. Seeing
attention this way demystifies it: it is nonparametric regression where the data
points are the other tokens, the thing a transformer layer does is smooth each
position's representation against the others, and ``what to attend to'' is just
``which keys are similar to this query.''

\begin{summarybox}
A language model is a distribution over sequences; the chain rule factorises it
into next-token conditionals, so the autoregressive model only needs
$p(x_t\mid x_{1:t-1})$ --- Markov-truncated ($n$-gram) or neural. Fitting it is
maximum likelihood, i.e.\ minimising cross-entropy, and perplexity is
$e^{\text{cross-entropy}}$ (Chapter~5). Word embeddings are not arbitrary:
word2vec implicitly factorises the shifted-PMI matrix, so semantics becomes
geometry (PCA for co-occurrence). Attention is the Nadaraya--Watson kernel
estimator with kernel $\exp(q^\top k/\sqrt d)$ --- a kernel-weighted average of
values, i.e.\ nonparametric regression over the other tokens.
\end{summarybox}

\end{document}
